% GENERATED FILE. Edit canonical lessons, then run npm run generate:books.
% canonical-source-hash: fnv1a64:b0955e009115201b
% canonical-lessons: RU-C130-bitter, RU-C130-salty, RU-C130-heavy, RU-C130-thin, RU-C130-fat

\chapter{Bitter, Salty, Heavy, Thin, Fat}
\label{ch:ru-bitter-salty-heavy-thin-fat}
\hlchaptermodality{130}

\begin{chapteropening}
\textbf{By the end of this chapter:} \emph{I can say bitter, salty, heavy, thin and fat.}

Complete the last lesson of chapter 130: I can say bitter, salty, heavy, thin and fat.
\end{chapteropening}

\section[\texorpdfstring{\ru{горький} (górkiy)}{горький (górkiy)}]{\ru{горький} (górkiy) — bitter}

\label{lesson:RU-C130-bitter}



\begin{quote}
Before the new one: say the Russian for wide, then the Russian for narrow.
\end{quote}

\subsection*{You'll want to know: \ru{горький}}
\textbf{\ru{горький}} (\emph{górkiy}) — \textquotedblleft{}bitter\textquotedblright{}.

\begin{cousinweb}
\textbf{\ru{горький}} goes with \emph{\ru{гореть}}, \textquotedblleft{}to burn\textquotedblright{} (Vasmer): a bitter taste burns. Vasmer compares Irish \emph{goirt}, \textquotedblleft{}bitter\textquotedblright{}.
\end{cousinweb}

\begin{grammarlens}[title={What you've built}]
One more word for this chapter.
\end{grammarlens}

\subsection*{Your turn}
\begin{itemize}
\raggedright
  \item \emph{Say it:} \emph{górkiy}
  \item \emph{Say it:} \emph{górkiy}, once more
  \item \emph{Say it:} say \emph{górkiy}
  \item \emph{Recall:} say the Russian for wide, then the Russian for narrow, then say \emph{górkiy} again
\end{itemize}

\begin{tcolorbox}[breakable,colback=teal!4,colframe=teal!35!black,title={Before you move on}]
What does \ru{горький} mean? (\textquotedblleft{}Bitter\textquotedblright{}.) Say it once more.
\end{tcolorbox}
\section[\texorpdfstring{\ru{солёный} (solyónyy)}{солёный (solyónyy)}]{\ru{солёный} (solyónyy) — salty}

\label{lesson:RU-C130-salty}



\begin{quote}
Before the new one: say the Russian for narrow, then the Russian for bitter.
\end{quote}

\subsection*{You'll want to know: \ru{солёный}}
\textbf{\ru{солёный}} (\emph{solyónyy}) — \textquotedblleft{}salty\textquotedblright{}.

\begin{cousinweb}
\textbf{\ru{солёный}} is from \emph{\ru{солить}}, \textquotedblleft{}to salt\textquotedblright{}, and \emph{\ru{соль}}, \textquotedblleft{}salt\textquotedblright{}, a cousin of Latin \emph{sāl} and English \emph{salt} (Vasmer).
\end{cousinweb}

\begin{grammarlens}[title={What you've built}]
One more word for this chapter.
\end{grammarlens}

\subsection*{Your turn}
\begin{itemize}
\raggedright
  \item \emph{Say it:} \emph{solyónyy}
  \item \emph{Say it:} \emph{solyónyy}, once more
  \item \emph{Say it:} say \emph{solyónyy}
  \item \emph{Recall:} say the Russian for narrow, then the Russian for bitter, then say \emph{solyónyy} again
\end{itemize}

\begin{tcolorbox}[breakable,colback=teal!4,colframe=teal!35!black,title={Before you move on}]
What does \ru{солёный} mean? (\textquotedblleft{}Salty\textquotedblright{}.) Say it once more.
\end{tcolorbox}
\section[\texorpdfstring{\ru{тяжёлый} (tyazhyólyy)}{тяжёлый (tyazhyólyy)}]{\ru{тяжёлый} (tyazhyólyy) — heavy}

\label{lesson:RU-C130-heavy}



\begin{quote}
Before the new one: say the Russian for bitter, then the Russian for salty.
\end{quote}

\subsection*{You'll want to know: \ru{тяжёлый}}
\textbf{\ru{тяжёлый}} (\emph{tyazhyólyy}) — \textquotedblleft{}heavy\textquotedblright{}.

\begin{cousinweb}
\textbf{\ru{тяжёлый}} $\leftarrow$ Proto-Slavic \emph{*tęžĭlŭ} — inherited. Vasmer relates it to Lithuanian \emph{tingùs}, \textquotedblleft{}lazy\textquotedblright{}, and Old Norse \emph{þungr}, \textquotedblleft{}heavy\textquotedblright{}.
\end{cousinweb}

\begin{grammarlens}[title={What you've built}]
One more word for this chapter.
\end{grammarlens}

\subsection*{Your turn}
\begin{itemize}
\raggedright
  \item \emph{Say it:} \emph{tyazhyólyy}
  \item \emph{Say it:} \emph{tyazhyólyy}, once more
  \item \emph{Say it:} say \emph{tyazhyólyy}
  \item \emph{Recall:} say the Russian for bitter, then the Russian for salty, then say \emph{tyazhyólyy} again
\end{itemize}

\begin{tcolorbox}[breakable,colback=teal!4,colframe=teal!35!black,title={Before you move on}]
What does \ru{тяжёлый} mean? (\textquotedblleft{}Heavy\textquotedblright{}.) Say it once more.
\end{tcolorbox}
\section[\texorpdfstring{\ru{худой} (khudóy)}{худой (khudóy)}]{\ru{худой} (khudóy) — thin}

\label{lesson:RU-C130-thin}



\begin{quote}
Before the new one: say the Russian for salty, then the Russian for heavy.
\end{quote}

\subsection*{You'll want to know: \ru{худой}}
\textbf{\ru{худой}} (\emph{khudóy}) — \textquotedblleft{}thin\textquotedblright{}.

\begin{cousinweb}
\textbf{\ru{худой}} $\leftarrow$ Proto-Slavic \emph{*xudŭ} — inherited; Old Russian \emph{xudŭ} meant \textquotedblleft{}bad, poor, small\textquotedblright{}. It is held to be related to Sanskrit \emph{kṣudrá-}, \textquotedblleft{}small, low\textquotedblright{} (Vasmer).
\end{cousinweb}

\begin{grammarlens}[title={What you've built}]
One more word for this chapter.
\end{grammarlens}

\subsection*{Your turn}
\begin{itemize}
\raggedright
  \item \emph{Say it:} \emph{khudóy}
  \item \emph{Say it:} \emph{khudóy}, once more
  \item \emph{Say it:} say \emph{khudóy}
  \item \emph{Recall:} say the Russian for salty, then the Russian for heavy, then say \emph{khudóy} again
\end{itemize}

\begin{tcolorbox}[breakable,colback=teal!4,colframe=teal!35!black,title={Before you move on}]
What does \ru{худой} mean? (\textquotedblleft{}Thin\textquotedblright{}.) Say it once more.
\end{tcolorbox}
\section[\texorpdfstring{\ru{толстый} (tólstyy)}{толстый (tólstyy)}]{\ru{толстый} (tólstyy) — fat, thick}

\label{lesson:RU-C130-fat}



\begin{quote}
Before the new one: say the Russian for heavy, then the Russian for thin.
\end{quote}

\subsection*{You'll want to know: \ru{толстый}}
\textbf{\ru{толстый}} (\emph{tólstyy}) — \textquotedblleft{}fat, thick\textquotedblright{}.

\begin{cousinweb}
\textbf{\ru{толстый}} $\leftarrow$ Proto-Slavic \emph{*tŭlstŭ} — inherited. Vasmer relates it to Lithuanian \emph{tulžti}, \textquotedblleft{}to swell\textquotedblright{}, and Latvian \emph{tùlzt}, \textquotedblleft{}to swell up\textquotedblright{}.
\end{cousinweb}

\begin{grammarlens}[title={What you've built}]
One more word for this chapter.
\end{grammarlens}

\subsection*{Your turn}
\begin{itemize}
\raggedright
  \item \emph{Say it:} \emph{tólstyy}
  \item \emph{Say it:} \emph{tólstyy}, once more
  \item \emph{Say it:} say \emph{tólstyy}
  \item \emph{Recall:} say the Russian for heavy, then the Russian for thin, then say \emph{tólstyy} again
\end{itemize}

\begin{tcolorbox}[breakable,colback=teal!4,colframe=teal!35!black,title={Before you move on}]
What does \ru{толстый} mean? (\textquotedblleft{}Fat, thick\textquotedblright{}.) Say it once more.
\end{tcolorbox}
